\section{Method}

\subsection{Change-Point Tests}

A change-point test takes a reported-results series ordered by time or batch
position and asks whether one or more locations split the sequence into regimes
with meaningfully different behavior. The feature may be a cumulative margin,
per-batch vote share, turnout rate, residual from a comparison model, or another
normalized RCOUNT-derived quantity. Candidate statistics include CUSUM scores,
segmentation costs, and likelihood-ratio change-point statistics.

The statistic is not interpreted alone. The report must record the ordering
rule, the feature set, the baseline population, and the operational note that
reporting order often creates benign changes. For example, mail ballots or
early-vote batches may be reported before election-day precincts; central-count
batches may arrive in blocks; small late batches may have different composition
than earlier high-volume batches. A flagged change point is therefore an
investigation item: it identifies the segment boundary and the observed shift,
then asks whether the reporting workflow explains it.

\subsection{Audit-Discrepancy Clustering}

A discrepancy-clustering test takes audit observations from the V-series side of
an RCOUNT package and asks whether discrepancies concentrate in operational
groups. The grouping variable may be scanner, batch, source file, tabulator,
audit board, or another unit preserved by the package.
The null is random scatter conditional on the relevant audit design and group
sizes. Candidate statistics include chi-square tests over grouped counts,
concentration indices, maximum-cell overrepresentation scores, and permutation
or rank-based clustering statistics.

The method consumes discrepancies as observations, not as certification
findings. It asks whether the pattern of discrepancies is more concentrated
than expected under the selected random-scatter baseline. A high score should
name the grouping unit, the observed concentration, the comparison population,
and the caveat. It should not say that the scanner, batch, or source caused the
problem, and it should not convert an audit discrepancy into proof of outcome
error or fraud.
